% ===================================================================
%  Cover letter — Elsevier Measurement (one page)
%  Compile: pdflatex cover_letter_measurement.tex
% ===================================================================
\documentclass[11pt,a4paper]{article}
\usepackage[utf8]{inputenc}
\usepackage[T1]{fontenc}
\usepackage{geometry}
\geometry{margin=2.2cm}
\usepackage{hyperref}
\pagestyle{empty}
\setlength{\parindent}{0pt}
\setlength{\parskip}{5pt}

\begin{document}

{\large\bfseries Cover Letter --- Measurement (Elsevier)}\hfill
{\small 24 September 2026}\\[2pt]
\hrule
\vspace{8pt}

Dawid Kucharski\\
Poznan University of Technology\\
Pozna\'n, Poland\\
dawid.kucharski@put.poznan.pl

\vspace{10pt}

Dear Editors,

I am pleased to submit the manuscript \textbf{``Quantum Limits of
Surface Metrology: Spatial Information Extraction by a Single Trapped
Ion''} for consideration as a research paper in \emph{Measurement}
(32-page main manuscript plus a 33-page Supplementary Material).

\textbf{Fit to the journal's scope.}  \emph{Measurement} invites
``novel achievements in all fields of measurement and instrumentation
science and technology''.  The manuscript proposes a new measurement
principle: uncertainty-budgeted reconstruction of conducting-surface
topography from the secular-frequency response of a single trapped ion.
It delivers the measurement-science core --- a defined measurand, an
explicit instrument transfer function $H_x(k;h)\propto k_x^2e^{-kh}$,
spatial-information limits, a topography--patch-potential
identifiability theorem, regularised inversion, a GUM-based uncertainty
budget, an SI traceability chain and five quantitative falsifiable
predictions --- formulated in the language of the VIM/GUM.  To our
knowledge, no previous trapped-ion surface-probing study has developed
such a combined metrological inverse framework.

\textbf{Rigour and validation.}  The transfer function is independently
cross-checked against exact Rayleigh--Floquet, boundary-element and
finite-difference electrostatic solvers (relative error $0.01\%$ at
$A=100$\,nm, growing as $0.14\,(kA)^2$), closed-form analytical
benchmarks and Monte Carlo uncertainty propagation ($95.1\%$ empirical
coverage under the stated input distributions), and is exercised on
published AFM topography data (reconstruction correlation $r=0.62$ at
$h=40$\,\textmu m improving to $r=0.78$ at $h=2$\,\textmu m).  This is
numerical evaluation, not experimental validation: the experimental
demonstration is transparently future work, specified through a costed
four-phase roadmap whose first phase runs on existing trapped-ion
hardware.

\textbf{Status.}  Single-author paper; not under consideration
elsewhere.  Code and data are openly available (GitHub repository
\url{https://github.com/dawidkucharski/qespm}; Zenodo deposit
\url{https://doi.org/10.5281/zenodo.22721000}; validation data from
Mendeley Data).

I would be glad to address any questions and thank you for considering
this work for \emph{Measurement}.

\vspace{10pt}

Sincerely,\\[12pt]

Dawid Kucharski\\
Poznan University of Technology, Pozna\'n, Poland

\end{document}
